\documentclass[10pt,a4paper]{amsart}
\usepackage[margin=19mm]{geometry}
\usepackage{amsmath,amssymb,mathtools}
\usepackage[scr=rsfs,cal=euler]{mathalfa}
\usepackage{xfrac}
\usepackage[unicode,hidelinks,bookmarks=false]{hyperref}
\newcommand{\ie}{i.e.\ }
\newcommand{\cf}{cf.\ }
\newcommand{\resp}{respectively\ }
\newcommand{\Subsection}{\subsection}
\providecommand{\nobreakdash}{\nobreak\hbox{-}}
\setlength{\emergencystretch}{2em}
\multlinegap0pt
\usepackage{amssymb}
\newtheorem{theorem}{Theorem}
\newcommand{\HS}{\mathrm{HS}}
\newcommand{\R}{\mathbb{R}}
\newcommand{\T}{\mathsf{T}}
\newcommand{\be}{\mathbf e}
\newcommand{\bw}{\mathbf w}
\newcommand{\norm}[1]{\lVert #1 \rVert}
\DeclareMathOperator{\tr}{tr}
\title{A Projection Identity for Simplices\\ Sharp Inequalities, Converse Results, and Affine Projections}
\author{Quang Hung Tran}
\date{Source: arXiv:2609.01226v1 (CC BY 4.0)}
\begin{document}
\maketitle

\noindent\emph{Source and licence.} A Projection Identity for Simplices\\ Sharp Inequalities, Converse Results, and Affine Projections. Version arXiv:2609.01226v1 (CC BY 4.0), distributed by arXiv under the Creative Commons Attribution 4.0 International licence, http://creativecommons.org/licenses/by/4.0/, as recorded in the arXiv licence metadata for this exact version and shown on the abstract page https://arxiv.org/abs/2609.01226v1. Full licence notice: \texttt{licenses/arxiv-2609.01226-CC-BY-4.0.txt}.\par\smallskip

\noindent\textit{Editorial scope.} Counted unit: the article's theorem thm:quantitative (source0.tex lines 570--588). The article's complete original proof of it is delivered with it (source0.tex lines 590--646, one \texttt{proof} environment of 57 lines). No mathematics is written, repaired or re-derived: the statement and the proof are the article's own lines, in the article's own notation, and no retained line is substituted or re-flowed.  Retained uncounted as the source-local context the proof needs: lines 503--519.  Nothing was omitted from any retained span, so the package carries no omission record.  No retained line contains a citation, so the wrapper carries no bibliography and no bibliography entry.  No retained line contains a figure environment, an image-inclusion command or drawing code, so no image is delivered and the package declares no asset dependency.  Editorial formatting only, in addition: an AMS \texttt{amsart} preamble that restates the article's own declarations and macros used by the retained lines, namely the environments \texttt{theorem} on separate counters, together with the standard packages those lines need.  The retained environments print in this document as Theorem 1 in order of appearance, whereas the article prints the same items with the numbering of its own full document.  The complete native source of the article as distributed in the arXiv e-print for this exact version is bundled unchanged as \texttt{sources/209086/source0.tex}.
\begin{theorem}\label{prop:howmany}
Fix $1\le k\le n-1$ and put
\[
D=\dim\operatorname{Sym}_n(\R)=\frac{n(n+1)}2.
\]
The least integer $N$ for which there exist $k$-dimensional subspaces
$L_1,\ldots,L_N$ with the following property is
\[
N=D-1=\frac{n(n+1)}2-1:
\]
whenever $\be_1,\ldots,\be_n$ are unit vectors and
\[
\sum_{i=1}^{n}\norm{Q_j\be_i}^2=k \text{ for } 1\le j\le N,
\]
where $Q_j$ denotes the orthogonal projection onto $L_j$, the vectors
$\be_1,\ldots,\be_n$ are an orthonormal basis.
\end{theorem}

\begin{theorem}\label{thm:quantitative}
Put
\[
c_n=
\begin{cases}
1/\sqrt n,& n\text{ even},\\[2pt]
1/\sqrt{n-1},& n\text{ odd}.
\end{cases}
\]
Then
\begin{equation}\label{eq:quantitative}
c_n\Theta
\le
\max_L|\Delta(L)|
\le
\sqrt{\frac{n-1}{n}}\,\Theta,
\end{equation}
where the maximum is over all lines through $S$. Both constants are best possible.
\end{theorem}

\begin{proof}
Put $X=M-I$. For $L=\R\bw$ with $\norm\bw=1$,
\[
\Delta(L)=\bw^{\T}X\bw,
\]
so
\[
\max_L|\Delta(L)|=\norm{X}_{\mathrm{op}}.
\]
Moreover,
\[
\norm{X}_{\HS}^2
=
\tr(M^2)-n
=
\tr(G^2)-n
=
\sum_{i\ne j}\cos^2\theta_{ij}
=
\Theta^2.
\]
Thus it remains to compare the operator and Hilbert--Schmidt norms of a nonzero symmetric
matrix with trace zero.

Let $x_1,\ldots,x_n$ be the eigenvalues of $X$ and put
$m=\max_i|x_i|$. For the upper bound in \eqref{eq:quantitative}, assume after changing
sign that $x_1=m$. Since $\sum_{i=2}^{n}x_i=-m$, Cauchy--Schwarz gives
\[
\sum_{i=2}^{n}x_i^2\ge\frac{m^2}{n-1}.
\]
Hence
\[
\norm{X}_{\HS}^2\ge\frac{n}{n-1}m^2,
\]
which is the required upper estimate. Equality occurs for the spectrum
\[
\Bigl(m,-\frac{m}{n-1},\ldots,-\frac{m}{n-1}\Bigr).
\]

For the lower bound, scale so that $m=1$. If $n$ is even, then simply
$\sum_i x_i^2\le n$, giving $m\ge\norm{X}_{\HS}/\sqrt n$; equality occurs when half of the
eigenvalues are $1$ and half are $-1$. If $n=2r+1$ is odd, write the positive and negative eigenvalues separately. One of the two
sign classes contains at most $r$ terms. Since $|x_i|\le1$ and the positive and negative
sums have the same absolute value, that common sum is at most $r$. Hence
\[
\sum_i x_i^2\le\sum_i|x_i|\le2r=n-1.
\]
This gives $m\ge\norm{X}_{\HS}/\sqrt{n-1}$, with equality for the spectrum consisting of
$r$ copies of $1$, $r$ copies of $-1$, and one $0$.

These equality patterns are realised by simplices as well. For each of
the trace-zero spectra above, choose a symmetric matrix $Y$ with that spectrum. For
sufficiently small $\varepsilon>0$, $I+\varepsilon Y$ is positive definite with trace $n$.
The same Schur--Horn argument used in Theorem~\ref{prop:howmany} produces unit edge vectors
with frame operator $I+\varepsilon Y$. Since both ratios in \eqref{eq:quantitative} are
scale invariant in $Y$, the constants are sharp.
\end{proof}

\end{document}
